\documentclass[24pt, a4paper, portrait]{article}

\usepackage[utf8]{inputenc}
\usepackage[a4paper,left=2cm,right=2cm,top=1cm,bottom=2cm]{geometry}
\usepackage{graphicx}
\usepackage{array}
\usepackage{subfig}

\graphicspath{{./img/}}

\begin{document}

\pagestyle{empty}

\raggedleft

\includegraphics[width=0.4\textwidth]{logo}

\vspace{1cm}
\sffamily
\centering
\huge

\textbf{Vor Betreten des Labors bitte registrieren!}

\vspace{0.5cm}
\Large
\textit{(VR4: Werkstatt)}
\vspace{0.5cm}

\raggedright

Um bei einem Covid-19-Fall in einem unserer Labore mögliche Infektionsketten zurückverfolgen zu können, müssen Beginn und Ende des Aufenthalts im Labor festgehalten werden.
Versuchsleiter*innen sind dafür verantwortlich, die Proband*innen ihrer Studien zu registrieren.
Beim Checkout muss anhand einer Checkliste bestätigt werden, dass der Arbeitsplatz vorschriftsgemäß desinfiziert wurde.

\medskip

Scannen Sie dazu den entsprechenden QR-Code oder folgen Sie diesen Links:

\textbf{wiki.mi.ur.de/lab/checkin}

\textbf{wiki.mi.ur.de/lab/checkin\_gast}

\textbf{wiki.mi.ur.de/lab/checkout}

\medskip

Füllen Sie das angezeigte Formular dann wahrheitsgemäß aus.
Es werden Name, Telephonnummer und E-Mail-Adresse, sowie der Aufenthaltszeitraum festgehalten.
Die erhobenen Daten werden nach spätestens 30 Tagen gelöscht.

\vspace{1cm}
\centering
%\includegraphics[width=0.5\textwidth]{qr/vr4_werkstatt_checkin}

\begin{figure}[!h]
    \centering
    \subfloat[Checkin]{\includegraphics[width=0.3\textwidth]{qr/vr4_werkstatt_checkin}}
    \hfill
    \subfloat[Checkin Gast]{\includegraphics[width=0.3\textwidth]{qr/vr4_werkstatt_checkin_gast}}
    \hfill
    \subfloat[Checkout]{\includegraphics[width=0.3\textwidth]{qr/vr4_werkstatt_checkout}}
\end{figure}

\vspace{1cm}
\centering
\huge
\textbf{Checkin und Checkout sind verpflichtend! \\ Ohne Registrierung ist das Betreten des Labors nicht gestattet!}

\end{document}
